\section{Introduction}

Dependent Object Types (DOT) is the calculus behind Scala's type system \cite{amin2016essence}. Its types are objects with type members and term members, and a type member $A$ of a variable $x$ can be selected as the type $x.A$. A binder $x : \dtyp{A}{S}{T}$ makes $S \sub x.A$ and $x.A \sub T$ available for whatever $S$ and $T$ the declaration names, and nothing forces $S \sub T$. Under a binder of type $\dtyp{A}{\top}{\bot}$ every type is a subtype of every other, and every term has every type.

This is the source of the difficulty in proving DOT sound. Such a binder can never be realized by a value, since no object has a member with those bounds, so a program never runs with it in scope. But the type system does not know that, and a proof by preservation and progress has to invert derivations that pass through subsumption at types like $x.A$. The published proofs deal with this in different ways. The WadlerFest proof characterizes the possible types of a value \cite{amin2016essence}. Rapoport et al.\ restrict contexts to inert ones and replace general typing by tight and invertible typing \cite{rapoport2017simple}. Rompf and Amin push transitivity back through derivations \cite{rompf2016soundness}, and Amin and Rompf move to a semantic model over a definitional interpreter \cite{amin2017definitional}. Each of these is an argument about the derivations of one particular rule set. Extending the calculus with paths \cite{rapoport2019path}, with mutation \cite{kabir2018kdot}, or with a new sort of member means extending the argument.

\paragraph{The route taken here.}
Compilers for languages with subtyping have a standard way of not inverting derivations: they make the derivations part of the program. System FC \cite{sulzmann2007fc} is the intermediate language of GHC, every use of a type equality in a source program becomes a coercion term of FC, and a small checker validates the output of every compiler pass. The metatheory of FC is the metatheory of a language without subsumption.

We do the same for DOT. A DOT typing derivation is translated into a term of \FCdot, a calculus in which every use of subtyping, of type equality, and of field presence is an explicit proof term. Terms of \FCdot type check with no subsumption rule. The proof terms erase to nothing, so the translated program computes what the source program computes.

What makes this more than a change of notation is the shape of the metatheory that results. Preservation and progress for \FCdot need one fact about evidence: over a typed store, closed evidence normalizes. A closed inclusion between two types normalizes to a head form, which is $\bot$ on the left, $\top$ on the right, a conversion between equal types, a function coercion, or an object coercion given by one entry per proposition of the target. A closed atom normalizes to a view, which lists the normal forms of the propositions its object type asserts about it. The proof is a structural induction on typing derivations of evidence. Evidence through bad bounds is not excluded. It normalizes like everything else, and when it is closed over a store it cannot exist, because no head form fits it. The impossibility of a closed $\top \le \bot$ is a corollary of the normalization theorem, proven once for all stores, and the translation never has to know that the contexts it works in are consistent.

\paragraph{Why not System FC.}
FC is a language of types without term dependence. DOT types mention variables, and the whole difficulty lives in that dependence. The design of \FCdot therefore departs from FC in three places. First, every term binder $x$ owns a block of type names $\sel{x}{\ell}$, one per label, and types refer to those names and to nothing else about terms. Second, an object type is a telescope of propositions over a self block: $\dtyp{A}{S}{T}$ becomes the object type asserting $\ple{S}{\sel{\self}{A}}$ and $\ple{\sel{\self}{A}}{T}$, and evidence between object types is a morphism that proves each target proposition from a source proposition. Third, evidence about the type of an atom can be eliminated at that atom: from $a : S$ and evidence for $S \le \obj{\tel}$ one obtains the propositions of $\tel$ at the block of $a$. The last rule is what a bad-bounds derivation becomes, and it is what a translation into existential packages in the style of F-ing modules \cite{rossberg2014fing} cannot express. Section~\ref{sec:counterexample} shows on a counterexample why opening a package at a fresh name is not enough.

\paragraph{Contributions.}
\begin{itemize}
\item \FCdot, an explicit-evidence calculus for DOT: types over binder blocks, telescopes of propositions as object types, evidence for inclusion, equality, and presence, template morphisms between object types, atoms, terms in monadic normal form, and a store machine (Section~\ref{sec:target}). The calculus has a decidable checker, proven complete.
\item The canonical forms theorem for closed evidence over a typed store, obtained from a structural normalizer, and the resulting preservation, progress, and erasure simulations (Section~\ref{sec:canonical}).
\item A translation of \DOTMNF, the WadlerFest calculus in monadic normal form, into \FCdot. It is a total function on derivations, homomorphic on types, type preserving, and coherent. Type safety of \DOTMNF and consistency of every reachable store are corollaries (Section~\ref{sec:translation}).
\item A mechanization in \Lean of everything above with no library dependencies, whose main theorems depend on propositional extensionality and quotient soundness only (Section~\ref{sec:mechanization}).
\end{itemize}

\paragraph{Scope.}
The source calculus has variable paths only, as in the WadlerFest presentation, and intersections are restricted to declaration-shaped types. Section~\ref{sec:conclusion} discusses what changes for paths.
